\section{Method}
\label{sec:method}
All models are frozen; we run inference only. Every claim below concerns the public V-JEPA ViT-L checkpoint unless stated otherwise.

\subsection{Surprise}
We reuse the surprise computation of Garrido et al.\ \cite{garrido2025}. A predictor receives the representations of $M$ context frames and predicts the representations of the next $N$ frames. The surprise at time $t$ is the $L_1$ distance between the predicted and observed representation, $s_t=\lVert \hat z_t - z_t\rVert_1$ (averaged over tokens). A clip's score aggregates the window surprises: the mean over windows for clean accuracy and invariance (H1), and the maximum for specificity and positive controls (H2, H3; event-aligned for H2). The default score takes, per window, the minimum over context lengths $C\in\{2,4,6,8,10\}$ frames. Ablations A1 (aggregation) and A2 (one fixed context length $C\in\{2,6,10\}$ at a time) are in Supplement Table~S5. Models: V-JEPA ViT-L (M1) \cite{garrido2025,bardes2024}, V-JEPA~2 ViT-L (M2) \cite{assran2025}, VideoMAE~V2 (M3, planned \cite{wang2023} but not run), M1 with random weights (M0r), and two baselines: frame-difference surprise $\overline{|f_t-f_{t-1}|}$ (B1) and the pixel error of copying the last frame (B2).

\subsection{Data}
The main evaluation uses the IntPhys (2019) dev split \cite{riochet2022}, 180 pairs from 90 scene quadruplets (30 per block O1--O3). Edit strengths were calibrated on the 60-video IntPhys~2 Debug split \cite{bordes2025}, which is not used for evaluation. IntPhys~2 Main is used descriptively, for clean accuracy only; no edited runs were made on it (no event-time annotation). The IntPhys (2019) data are used as released for evaluation, and no videos are redistributed; the IntPhys~2 dataset card gives CC BY-NC 4.0. IntPhys~2 is used for evaluation only. In IntPhys dev, each impossible video is built by splicing frames of its possible partner videos (established by pixel hashing of the dev videos: 120 single and 60 double splices, none unspliced). This describes the dataset and is separate from the hypothesis, in the Discussion, that M1 relies on the splices.

\subsection{Edits}
Each edit is a deterministic function of frames, parameters and seed, applied before model preprocessing at three strengths (low, mid, high). Blur and low illumination at several graded severities are a common way to test vision models under degradation \cite{koushik2025,randieri2025}. Blur, JPEG and noise strengths are set on Debug frames to SSIM $0.95/0.85/0.70$ (JPEG high, $q=3$, reaches $0.689$ because quality is an integer); brightness gain to a mean CIELAB $L^*$ change of $5/10/20$; centre-crop scale is fixed at $0.9/0.8/0.7$. Horizontal flip, frame-rate halving and the dither floor have a single setting. \emph{Invariance edits} are applied with identical parameters and seed to both videos of a pair: brightness/colour shift, Gaussian blur, JPEG re-compression, additive Gaussian noise, centre crop and resize, horizontal flip, frame-rate halving (degenerate, see item~(f) of the deviations below). \emph{Abrupt edits} are applied to possible videos only, at an event time $t^*$: a persistent lighting step, camera jolt, or sudden zoom. The event time $t^*$ is the first frame at which the mean absolute pixel difference $D(t)$ between the impossible video and its prefix-sharing possible partner exceeds $\max(2/255,\ \text{the 99.9th percentile of } D \text{ over the first 10 frames of all pairs})$. If $D$ never exceeds this floor, the rule falls back to a seeded draw from the middle half of the clip. \emph{Positive controls}: vertical flip and time reversal of possible clips.

\emph{Magnitude matching.} For each pair, the violation's pixel change is $|\Delta|_{\mathrm{viol}}=\overline{|v_{\mathrm{imp}}-v_{\mathrm{pos}}|}$ over frames after the event, and the edit strength is solved so that the edit's mean pixel change over the same window equals the pair's violation magnitude (closed form for lighting, bisection for jolt and zoom); mid strength is $1\times$ this value, low $0.5\times$, high $2\times$. The pre-registered target is a median achieved-to-target ratio in $[0.95,1.05]$ with at least 90\% of pairs in $[0.9,1.1]$. Matching uses dev pixels only, never model output.

\subsection{Metrics and statistics}
\emph{Pairwise accuracy} is the fraction of pairs with $S(\text{impossible})>S(\text{possible})$; \emph{flip rate} is the fraction whose verdict differs from the clean one. We use McNemar tests against clean with Holm correction across edits. \emph{Specificity AUROC} is $P(\Delta_{\mathrm{viol}}>\Delta_{\mathrm{abrupt}})$ with $\Delta_{\mathrm{viol}}=S(\text{imp})-S(\text{pos})$ and $\Delta_{\mathrm{abrupt}}=S(\text{pos+edit})-S(\text{pos})$, with a 10{,}000-resample bootstrap 95\% CI; the $D_1$ and $D_2$ contrasts use a one-sided paired bootstrap $p$-value (floor $2\times10^{-4}$) and a paired CI. \emph{Detection rate} for positive controls is the fraction of videos with $S(\text{edited})>S(\text{original})$ with a binomial CI. Edit randomness uses 3 seeds (sign, colour direction, noise realisation, jolt direction); deterministic edits run once. The bootstrap resamples scene quadruplets, stratified by block. Ties are scored 0.5 (``tie-half'') when two scores differ by at most $2\times10^{-4}$ times the median clean score of that type; a strict convention ($<$, as in the released code) and an exact-tie convention ($\tau=0$) are sensitivity analyses, and we report the tie-half and strict figures for the dither floor.

\subsection{Pre-registration}
Edit parameters, seeds, metrics, analysis plan and thresholds were frozen on 7 October 2026 at commit \texttt{ee63aee} (tag \texttt{prereg-frozen}), before any edited-condition output existed. Strengths were calibrated on IntPhys~2 Debug only. Decision rules, as frozen:
\begin{itemize}
\item \emph{Invariance (H1).} An edit is ``violated'' if accuracy drops by at least 5 points with a Holm-adjusted scene-cluster bootstrap $p<0.05$, or if its flip rate is at least 10\% and at least 5 points above the dither floor with Holm-adjusted $p<0.05$; it is ``invariant'' if the accuracy lower bound exceeds $-5$ points and the upper bound of flip minus floor is below 5 points. The Holm family has 14 tests (7 edits $\times$ accuracy and flip). If the dither floor is at least 10\%, H1 for that model is descriptive.
\item \emph{Specificity (H2).} Let $A$ be the specificity AUROC of M1, and $A_0$ the clean pairwise accuracy of M1 under the same event-aligned max aggregation (detection precondition). Define $D_1=A_{\mathrm{M1}}-A_{\mathrm{B1}}$ and $D_2=A_{\mathrm{M1}}-A_{\mathrm{B2}}$. ``Specific'' requires $A\ge0.80$, both $D_1$ and $D_2$ significant under Holm correction over these two tests, and per-type AUROC of at least 0.65; $0.65<A<0.80$ is partial; $A\le0.65$ is not specific. If $A_0<0.85$ and the table yields not specific or partial, the outcome is ``detection-limited'': specificity cannot be assessed and no claim is made in either direction. This Holm-2 rule replaced an earlier intersection-union rule before the freeze.
\item \emph{Positive controls (H3, gate G2).} Vertical flip must raise M1's surprise in at least 80\% of possible videos, with the lower confidence bound against horizontal flip above 0.5.
\end{itemize}
A model is H1-eligible only if its clean accuracy is at least 80.0\%. Analyses outside this plan are labelled exploratory.

\subsection{Deviations, amendments and known design weaknesses}
\label{sec:deviations}
Each item departs from the plan or limits what a result can show, and is carried into Section~\ref{sec:discussion}. Full descriptions are in Supplement Section~S-I.
(a)~\emph{Checkpoint gap (G1$'$).} The original gate (at least 90\% on IntPhys dev) was not met: the public ViT-L checkpoint reaches 81.7\% (ViT-H 88.3\%, tie-half under the pre-registered threshold). The checkpoints behind the published results are not released and the gap is unexplained; we do not claim a reproduction of the published 98\%, which is for ViT-H with per-property context length. The replacement gate G1$'$ checks the pipeline only, and all three of its checks passed.
(b)~\emph{Positive controls (G2 not met).} Vertical flip raised M1's surprise in 21.1\% of possible videos against the 80\% gate, and pairwise accuracy on flipped pairs is not lower than clean. Three readings remain open (the edits are weak violations for this model; the pipeline is insensitive; the signal acts as a splice-discontinuity detector), so M1's sensitivity to physical impossibility is not established. By human decision (Q10-A, logged 2026-10-08) the study continued with G2 not met, so every H1 and H2 statement carries this caveat. The random-weight network M0r responds in 87.8\% of videos where about 50\% was expected, so M1 minus M0r is not a measure of learned sensitivity.
(c)~\emph{Magnitude matching in H2.} The frozen event rule fell back to a seeded random onset in 178 of 180 pairs. This is a failure of the rule, not a finding that no event exists. The median matched magnitude is 0.33 of 255 grey levels, so the edits are near-threshold; for global lighting, 84 of 180 entries are no-ops, so $A$ is reported with and without that type. Both pixel baselines fall below 0.5 for a reason we did not verify.
(d)~\emph{M2 wrapper.} The first V-JEPA~2 run used mask token index~1, but only token~0 is trained. After the freeze we changed it to~0 (human-approved, logged before any edited M2 output was opened) and reran the clean split: 46.9\% before, 60.8\% after. Both are reported; M2 is not H1-eligible.
(e)~\emph{Invariance is descriptive.} The dither floor of 21.5\% (15.6\% strict) exceeds 10\%, so every edit is read through its accuracy change and its flip rate relative to the floor. The floor was measured in fp16 only. For JPEG the verdict depends on the tie convention, and for zoom a retention rule leaves 112 of 180 pairs.
(f)~\emph{Frame-rate edit (design error).} The pre-registered edit reproduces the skip-4 clean reference bit-identically, so its zero change cannot be read as invariance; it is excluded from all counts.
(g)~\emph{Other.}
\begin{itemize}
\item Holm families. H1 primary: 14 tests. H1 dose-response: 30 tests for M1 (the pre-registration names 21, but only 15 graded cells exist) and 22 for M2 (11 cells, two tests each). H2 dose-response: 24 tests at cell level and 6 for the pooled contrasts. The M2 dose cells are seed 1 only.
\item Precision. H1 eligibility was judged on the fp32 clean run and contrasts use fp16; the G1$'$ ladder compares ViT-L in fp32 with ViT-H in fp16.
\item Logistics. Kernel bundles were rebuilt at later commits with the commit allowlist extended (human decision Q10-C: ``proceeds''); the per-pair matching file was generated after the freeze from dev pixels only; the Colab fallback was not used. H2 was run as frozen, with its weaknesses disclosed, by human decision (Q11).
\end{itemize}
Exploratory analyses are labelled wherever they appear; the IntPhys~2 Main clean run is descriptive, as pre-registered.
